\documentclass{article}
\usepackage[UTF8]{ctex}
\usepackage{amsmath,amssymb,amsfonts}
\usepackage{enumitem}
\usepackage{geometry}
\geometry{a4paper, margin=1in}
% === 错题本排版优化（不参与题目标识） ===
\setlength{\parindent}{0pt}
\setlength{\parskip}{0.35em plus 0.10em minus 0.08em}
\setlength{\emergencystretch}{3em}
\setlist{itemsep=0.25em, topsep=0.40em, parsep=0pt}
\allowdisplaybreaks[2]
% === 错题本排版优化结束 ===
\begin{document}
\section*{同济大学2025年高等代数真题}
\section*{35．同济大学2025年高等代数试卷}
一、填空题．每题5分，共50分．

\begin{itemize}

\subsection{题目1}
  \item[1．] 设 \(m(x)\) 为 \(n\) 阶复方阵 \(A\) 的极小多项式，且 \(m(x)=f(x) g(x)\) ，其中最大公因式 \((f(x), g(x))=1\)，那么 \(R(f(A))+R(g(A))=\) \(\_\_\_\_\)．

\subsection{题目2}
  \item[2．] 设方阵 \(A=\left(\begin{array}{cc}-1 & 0 \\ 2 & -1\end{array}\right)\) ，那么 \(A^{2024}=\) \(\_\_\_\_\)．
\end{itemize}

\subsection{题目3}
3．设 \(A\) 是 6 阶方阵，\(B=\left(\begin{array}{ccc}2 A & O & 2 A \\ 2 A & A & 4 A \\ O & 2 A & 5 A\end{array}\right)\) ，如果 \(|A|=2\) ，那么 \(|B|=\)

\subsection{题目4}
4．若 \(A\) 和 \(B\) 是 3 阶复方阵，且 \((A B)^{3}=O\) ，那么 \(R\left((B A)^{2}\right)\) 所有可能的取值为 \(\_\_\_\_\)．

\subsection{题目5}
5．设 \(A\) 是 \(n\) 阶方阵且 \(R(A)=1, \operatorname{Tr}(A)=0\) ，记 \(E_{n}\) 为 \(n\) 阶单位阵，那么 \(E_{n}+A\) 是否可逆 \(\_\_\_\_\)．

\subsection{题目6}
6．给定 \(\mathbb{R}^{4}\) 的如下两个子空间

\[
\begin{gathered}
W_{1}=\left\{\left(x_{1}, x_{2}, x_{3}, x_{4}\right)^{\mathrm{T}} \mid x_{1}+x_{2}-x_{4}=0\right\} ; \\
W_{2}=\left\{\left(x_{1}, x_{2}, x_{3}, x_{4}\right)^{\mathrm{T}} \mid 3 x_{1}-x_{2}+2 x_{3}=0,5 x_{1}+x_{2}-3 x_{4}=0\right\} .
\end{gathered}
\]

则 \(\operatorname{dim}\left(W_{1}+W_{2}\right)=\) \(\_\_\_\_\)．

\subsection{题目7}
7．设矩阵

\[
A=\left(\begin{array}{ccc}
1 & -1 & 1 \\
1 & 1 & 0 \\
1 & 1 & 1
\end{array}\right)
\]

存在正交阵 \(Q\) 和对角元都是正数的上三角阵 \(R\) ，使得 \(A=Q R\) ，则 \(Q=\) \(\_\_\_\_\)．

\subsection{题目8}
8．设 \((V,())\) 是一个三维欧氏空间，\(e_{1}, e_{2}, e_{3}\) 为 \(V\) 的一组基，且内积在 \(\left(e_{1}, e_{2}, e_{3}\right)\) 下的度量矩阵为

\[
A=\left(\begin{array}{lll}
2 & 1 & 1 \\
1 & 2 & 1 \\
1 & 1 & 4
\end{array}\right)
\]

则当 \(x^{2}+y^{2}+z^{2}=1\) 时，向量 \(x e_{1}+y e_{2}+z e_{3}\) 的长度最大值为 \(\_\_\_\_\)．

\subsection{题目9}
9．设矩阵 \(A=\left(\begin{array}{ccc}1 & 1 & 1 \\ a & b & c \\ a^{2} & b^{2} & c^{2}\end{array}\right)\) ，其中 \(a, b, c\) 是实数，且 \(a>b>c\) ，令 \(B=A^{\mathrm{T}} A\) ，有下面几个说法：\\
（1）B 的特征值一定是实数．\\
（2）B 一定可对角化．\\
（3）一定成立 \(R(A)=R(B)\) ．\\
（4）二次型 \(X^{\mathrm{T}} B X\) 一定正定。上面说法正确的有 \(\_\_\_\_\)。

\subsection{题目10}
10．对 \(n\) 阶复方阵 \(A\) ，定义复线性空间 \(C(A)=\left\{X \in M_{n}(\mathbb{C}) \mid X A=A X\right\} \subseteq M_{n}(\mathbb{C}), M_{n}(\mathbb{C})\) 是所有 \(n\) 阶复方阵构成的线性空间。若 \(A\) 的 Jordan 标准型中属于不同特征值的 Jordan 块只有一个，那么 \(C(A)\) 的维数能取到的所有可能值为 \(\_\_\_\_\)．

二．解答题。\\

\subsection{题目11}
1．（20 分）设矩阵 \(A=\left(\begin{array}{cccc}1 & 1 & 1 & -2 \\ 1 & 2 & 3 & -5 \\ 2 & 1 & \lambda & -1 \\ 1 & 0 & -1 & \mu\end{array}\right)\) ，向量 \(\beta=(1,0,3,2)^{\mathrm{T}}\) ，其中 \(\lambda\) 和 \(\mu\) 是实数。设 \(\alpha_{1}, \alpha_{2}, \alpha_{3}, \alpha_{4}\)是 \(A X=\beta\) 的互不相同的解，\(X=\left(x_{1}, x_{2}, x_{3}, x_{4}\right)^{\mathrm{T}}\) ．

\begin{itemize}
  \item[（1）] 求 \(\lambda\) 和 \(\mu\) 的值，并写出 \(A X=\beta\) 的通解。
  \item[（2）] 求列满秩阵 \(H\) 和行满秩阵 \(L\) ，使得 \(H\) 是 \(A\) 的子阵，且 \(A=H L\) 。
  \item[（3）] 证明：\(\alpha_{1}-\alpha_{2}, \alpha_{2}-\alpha_{3}, \alpha_{3}-\alpha_{4}\) 线性相关。
\end{itemize}

\subsection{题目12}
2．（20 分）设实二次型

\[
f\left(x_{1}, x_{2}, x_{3}\right)=x_{1}^{2}+x_{2}^{2}+x_{3}^{2}-2 x_{1} x_{2}-2 a x_{1} x_{3}-2 x_{2} x_{3}
\]

在单位球面 \(\left\{\left(x_{1}, x_{2}, x_{3}\right)^{\mathrm{T}} \in \mathbb{R}^{3} \mid x_{1}^{2}+x_{2}^{2}+x_{3}^{2}=1\right\}\) 上取值最大和最小为2和－1，\(a\) 是正实数。

\begin{itemize}
  \item[（1）] 求 \(a\) 的值．
  \item[（2）] 求正交变换 \(\left(x_{1}, x_{2}, x_{3}\right)^{\mathrm{T}}=P\left(y_{1}, y_{2}, y_{3}\right)^{\mathrm{T}}\) ，将二次型 \(f\) 化为标准形，并写出标准形．
\end{itemize}

\subsection{题目13}
3．（20 分）设 \(n \geq 3\) 为整数，\(n\) 阶矩阵

\[
A=\left(\begin{array}{ccccc}
0 & & & & -a_{0} \\
1 & 0 & & & -a_{1} \\
& 1 & 0 & & -a_{2} \\
& & & \ddots & \vdots \\
& & & 1 & 0-a_{n-2} \\
& & & & 1-a_{n-1}
\end{array}\right)
\]

这里 \(a_{0}, a_{1}, \cdots, a_{n-1}\) 为有理数．证明：

\begin{itemize}
  \item[（1）] 若多项式 \(f(x)=x^{n}+a_{n-1} x^{n-1}+\cdots+a_{1} x+a_{0}\) 在 \(\mathbb{Q}\) 上不可约，作为复数域，\(A\) 可相似对角化。
  \item[（2）] 若对 \(i=0,1, \cdots, n-1\) ，有 \(a_{i}=\frac{n!}{i!(n-i)!}(-a)^{n-i}\) ，其中 \(a\) 是一个复数，求 \(A\) 的 Jordan 标准形．
\end{itemize}

\subsection{题目14}
4．（15 分）设 \(\varphi\) 是线性空间 \(V\) 上的线性变换，如果 \(\varphi\) 的不变子空间只有 \(\{0\}\) 和 \(V\) ，称 \(\varphi\) 是不可约的．若 \(V\) 不能写成 \(\varphi\) 的两个非零子空间的直和，则称 \(\varphi\) 不可分解。若 \(\varphi\) 的每个不变子空间都有不变补空间，则称 \(\varphi\) 是半单的．设 \(\varphi\) 是半单的，\(W\) 是 \(\varphi\) 的非零不变子空间．证明：

\begin{itemize}
  \item[（1）] \(\varphi\) 在 \(W\) 上的限制 \(\varphi \mid W: W \rightarrow W\) 也是半单的。
  \item[（2）] 线性变换 \(\varphi \mid W\) 不可分解当且仅当 \(\varphi \mid W\) 不可约。
\end{itemize}

\subsection{题目15}
5．（15 分）设 \(V\) 是 \(n\) 为欧氏空间，\(\alpha_{1}, \alpha_{2}, \cdots, \alpha_{n} \in V\) 。证明：\(\alpha_{1}, \alpha_{2}, \cdots, \alpha_{n}\) 是 \(V\) 的标准正交基的充要条件是对任意的 \(\alpha \in V\) ，有 \(\|\alpha\|^{2}=\sum_{i=1}^{n}\left(\alpha, \alpha_{i}\right)^{2}\) ．

\subsection{题目16}
6．（10 分）设整数 \(n \geq 2, A\) 是 \(n\) 阶实矩阵，且存在 \(a \in[-2,2]\) ，使得 \(A^{2}+a A+E_{n}=O\) ．证明：对任意的正整数 \(m \geq 2\) ，存在唯一的 \(a_{m} \in[-2,2]\) ，使得 \(A^{2 m}+a_{m} A^{m}+E_{n}=O\) 。
\end{document}
